% ---
% file: docs/latex/constantes/constantes.tex
% description: Catálogo de constantes matemáticas organizado por nivel didáctico y propiedades de irracionalidad/trascendencia.
% type: doc/manual
% version: 1.0.0
% date: 2026-08-24
% covers:
%   - src/data/constants.py
% relations:
%   - docs/fisica/constantes/constantes.md
%   - docs/ARCHITECTURE.md
% keywords:
%   - constantes-matematicas
%   - numeros-irracionales
%   - numeros-trascendentes
%   - constantes-geometricas
% ---

\documentclass[11pt,a4paper]{article}
\usepackage[utf8]{inputenc}
\usepackage[T1]{fontenc}
\usepackage[spanish,es-nodecimaldot,es-noquoting]{babel}
\usepackage{amsmath,amssymb,amsfonts}

% Configuración de márgenes y espaciado estándar
\setlength{\textwidth}{16.5cm}
\setlength{\textheight}{24cm}
\setlength{\oddsidemargin}{-0.25cm}
\setlength{\evensidemargin}{-0.25cm}
\setlength{\topmargin}{-1cm}
\setlength{\parindent}{0pt}
\setlength{\parskip}{5pt}

\begin{document}

\section*{Catalogo de constantes matematicas}
\addcontentsline{toc}{section}{Catalogo de constantes matematicas}

Las constantes exactas se muestran mediante una definicion; los decimales son aproximaciones truncadas o redondeadas. La clasificacion distingue lo demostrado de lo que permanece abierto.

\section*{1. Nivel basico: geometria, aritmetica y calculo elemental}
\addcontentsline{toc}{section}{1. Nivel basico: geometria, aritmetica y calculo elemental}

\begin{itemize}
\footnotesize
  \item \textbf{Pi ($\pi$):} $C/d=3.14159265358979323846264338327950\ldots$; real, irracional y trascendente.
  \item \textbf{Numero aureo ($\varphi$):} $(1+\sqrt5)/2=1.61803398874989484820458683436563\ldots$; algebraico de grado 2.
  \item \textbf{Proporcion de plata ($\delta_S$):} $1+\sqrt2=2.41421356237309504880168872420969\ldots$; algebraico de grado 2.
  \item \textbf{Constante de Pitagoras ($\sqrt2$):} diagonal de un cuadrado unitario, $x^2=2$;\newline $1.41421356237309504880168872420969\ldots$.
  \item \textbf{Constante de Teodoro ($\sqrt3$):} diagonal de un rectangulo de lados $1$ y $\sqrt2$ (o diagonal espacial de un cubo unitario), $x^2=3$;\newline $1.73205080756887729352744634150587\ldots$.
  \item \textbf{Base del logaritmo natural ($e$):} $\lim_{n\to\infty}(1+1/n)^n=\sum_{k=0}^\infty1/k!=\newline 2.71828182845904523536028747135266\ldots$; trascendente.
\end{itemize}

\section*{2. Nivel intermedio: analisis, numeros y combinatoria}
\addcontentsline{toc}{section}{2. Nivel intermedio: analisis, numeros y combinatoria}

\begin{itemize}
  \item \textbf{Euler-Mascheroni ($\gamma$):} $\lim_{n\to\infty}(H_n-\ln n)=0.57721566490153286060651209008240\ldots$; su irracionalidad aun no esta demostrada.
  \item \textbf{Catalan ($G$):} $\beta(2)=\sum_{n=0}^\infty(-1)^n/(2n+1)^2=0.91596559417721901505460351493238\ldots$; irracionalidad no resuelta.
  \item \textbf{Apery ($\zeta(3)$):} $\sum_{n=1}^\infty n^{-3}=1.20205690315959428539973816151144\ldots$; irracional, trascendencia no demostrada.
  \item \textbf{Liouville ($c$):} $\sum_{n=1}^\infty 10^{-n!}=0.110001000000000000000001\ldots$; el primer numero trascendente construido explicitamente.
  \item \textbf{Gompertz ($\delta_G$):} $\int_0^\infty e^{-t}/(1+t)\,dt=-e\operatorname{Ei}(-1)=0.59634736232319407434107849936927\ldots$.
  \item \textbf{Glaisher-Kinkelin ($A$):} constante definida por la expansion asintotica del hiperfactorial, $1.28242712910062263687534256886979\ldots$; no se debe afirmar que su trascendencia esta demostrada.
  \item \textbf{Khinchin ($K_0$):} media geometrica de los cocientes de fracciones continuas para casi todo real, $2.68545200106530644530971483548179\ldots$; no se conoce su irracionalidad.
  \item \textbf{Meissel-Mertens ($M$):} $\lim_{x\to\infty}(\sum_{p\le x}1/p-\ln\ln x)=\newline 0.26149721284764278375542683860869\ldots$.
  \item \textbf{Brun para primos gemelos ($B_2$):} suma de los inversos de ambos primos de cada pareja gemela, aproximadamente $1.902160583104$; la cota depende del calculo numerico y no es un ``rango analitico'' de la constante.
\end{itemize}

\section*{3. Nivel avanzado: dinamica, algoritmos e informacion}
\addcontentsline{toc}{section}{3. Nivel avanzado: dinamica, algoritmos e informacion}

\begin{itemize}
  \item \textbf{Feigenbaum ($\delta$):} razon limite de intervalos sucesivos de duplicacion de periodo,\newline $4.66920160910299067185320382046620\ldots$; universal para una clase de mapas, no ``analiticamente exacta'' en forma cerrada.
  \item \textbf{Feigenbaum ($\alpha$):} razon de escalas espaciales de atractores sucesivos,\newline $2.50290787509589282228390287321821\ldots$; la convencion puede cambiar el signo.
  \item \textbf{Conway ($\lambda$):} crecimiento de \textit{look-and-say}, $L_n\sim C\lambda^n$,\newline $1.30357726903429639125709911215255\ldots$; algebraico de grado 71.
  \item \textbf{Golomb-Dickman ($\lambda_{GD}$):} fraccion esperada del ciclo mas largo de una permutacion aleatoria,\newline $0.62432998854355087099293638310083\ldots$.
  \item \textbf{Chaitin ($\Omega_U$):} $\sum_{p\in\operatorname{dom}(U)}2^{-|p|}$, dependiente de la maquina universal $U$, no computable y aleatoria algoritmicamente. Un decimal atribuido a una ``maquina estandar'' no es universal y debe citar su construccion concreta.
\end{itemize}

\subsection*{Cobertura didactica}
\addcontentsline{toc}{subsection}{Cobertura didactica}

\begin{itemize}
  \item \textbf{Basico:} constantes geometricas y $e$, adecuadas para aritmetica, algebra, trigonometria y calculo inicial.
  \item \textbf{Intermedio:} $\gamma$, Catalan, $\zeta(3)$ y constantes de teoria de numeros, con series y limites convergentes.
  \item \textbf{Avanzado:} Feigenbaum, Conway, Golomb-Dickman y Chaitin, que requieren dinamica, asintotica, combinatoria o teoria de la computacion.
\end{itemize}


\end{document}
